
\documentclass[12pt,a4paper,heading]{report}

\usepackage[brazil]{babel}
\usepackage{ucs}
\usepackage[utf8x]{inputenc}
\usepackage[perpage]{footmisc}
\usepackage{float,epsfig,amsmath,amsfonts,latexsym,amstext,amsbsy,amsthm}
\usepackage[languagenames,fixlanguage]{babelbib}
\usepackage{color,bbm,tocbibind,setspace,url,ifpdf,graphicx,fancyhdr,fancyvrb}
\usepackage[a4paper,body={21.0cm,29.7cm},left=2.0cm,right=1.0cm,top=2.0cm,bottom=2.0cm]{geometry}
\usepackage{soul}
\usepackage{wrapfig}
\usepackage{hyperref}

\pdfinfo{
   /CreationDate (D:20111019170432)
}

\definecolor{darkblue}{rgb}{0.3,0.3,0.3}
\hypersetup{
	pdfcreator={Apostilator v0.3-20120211},
	pdfauthor={Reinaldo Gil Lima de Carvalho},
	pdftitle={Ambiente de Desenvolvimento JEE},
	pdfsubject={Ambiente de Desenvolvimento JEE},
	pdfkeywords={Java, JEE, Eclipse, Maven, Framework Demoiselle, JBoss, GlassFish, Application Server},
	plainpages=true,
	bookmarksopen=true,
	bookmarksopenlevel=1,
	pdfstartview=FitV,
	bookmarksnumbered=true,
	linktocpage=true,
	colorlinks=true,
	urlcolor=darkblue,
	linkcolor=darkblue,
	filecolor=darkblue,
	citecolor=darkblue
}

\pdfpagewidth=\paperwidth
\pdfpageheight=\paperheight

\newcounter{SteP}
\setcounter{SteP}{1}

%\pagecolor{grey}
%\color{red} % altera a cor do texto

\definecolor{myColor}{rgb}{0.9,0.9,0.9}
\sethlcolor{myColor}

\DefineVerbatimEnvironment%
  {VerbatimNumerado}{Verbatim}
    {numbers=left,numbersep=5pt,
       frame=single,framerule=0.3mm}

\DefineVerbatimEnvironment%
  {BoxVerbatim}{Verbatim}
    {numbers=none,frame=single}

% \DefineVerbatimEnvironment%
%  {MyVerbatim3}{Verbatim}
%    {numbers=none,frame=single,backgroundcolor=myColor}
      
\fancypagestyle{plain}{
\fancyhf{}
\renewcommand{\headrulewidth}{0pt}
\fancyfoot[R]{\thepage}
\fancyfoot[L]{\bf \small \slshape Ambiente de Desenvolvimento JEE} %% @chapterfooter
}

\fancypagestyle{rcsfooters}
{
    \fancyhf{}
    \renewcommand{\headrulewidth}{0.5mm}
    \renewcommand{\footrulewidth}{0.5mm}
    \renewcommand{\headwidth}{\textwidth}
    \renewcommand{\linewidth}{\textwidth}

    %
    %\fancyhead[L]{\bf Curso III: Emacs}
    %\fancyhead[R]{\includegraphics[scale=0.25]{figs/capa.eps}}
    \fancyhead[L]{\bf Java Enterprise} %% @header
    \fancyhead[R]{\bf \slshape \nouppercase{\leftmark}}

    \fancyfoot[R]{\thepage}
    \fancyfoot[C]{}
    \fancyfoot[L]{\bf \small \slshape Reinaldo de Carvalho - reinaldoc@gmail.com} %% @footer
}

\newcommand{\lpi}{\begin{wrapfigure}{l}{1.5cm}\centering \includegraphics[width=1.3cm]{titlepage-imgs/lpi.png} \end{wrapfigure}}
\newcommand{\rh}{\begin{wrapfigure}{l}{1.5cm}\centering \includegraphics[width=1.3cm]{titlepage-imgs/redhat.png} \end{wrapfigure}}

\PrerenderUnicode {^^c3^^b3}

\begin{document}\pagestyle{rcsfooters}

\thispagestyle{empty}
\thispagestyle{empty}

\begin{center}
  \begin{figure}[H]
    \begin{center}
      \hspace{-1cm}
      \includegraphics[width=0.55\columnwidth]{titlepage}
  ~\cite{FrontPageIMG}
    \end{center}      
  \end{figure}

  \vskip 2cm

  \hspace{-1cm}
  \begin{minipage}[c]{17cm}
    \begin{center}

{\huge \sf {Tecnologias de desenvolvimento em Java}\vskip 0.15cm %% @title

{\large Reinaldo Gil Lima de Carvalho} %% @author
}
    \end{center}
  \end{minipage}

  \vskip 14cm

  {\huge Junho 2012} %% @date

\end{center}

\newpage
\tableofcontents
\newpage
\newpage \chapter{Introdução ao curso}\setcounter{SteP}{1}

    Seja bem-vindo ao {\bf Curso de Drupal}. Este curso esta sendo
fomentado pelo {\it Xxxxx}
da {\it Zzzzz} e ministrado por {\it Reinaldo Gil Lima
de Carvalho}. Realiza(ou)-se de {\it xx/yy/2017 a xx/yy/2017}.

    Os procedimentos descritos neste material foram validados sob a
distribuição GNU/Linux Debian Stretch~\cite{Debian}, todavia a base teórica ministrada é o
conhecimento fundamental para a aplicabilidade desses procedimentos sob
qualquer sistema operacional. Busca-se que o participante esteja apto a
utilizar o conhecimento adquirido no ambiente que lhe for mais adequado.

    Neste capítulo, serão abordados os seguintes temas: licenciamento deste
material, origem do {\it software} utilizado (executável ou do código fonte)
e configurações essenciais da distribuição Debian GNU/Linux.

\section{Licença do material}\setcounter{SteP}{1}

    Todas as marcas registradas são de propriedade de seus respectivos
detentores, sendo apenas citadas neste material educacional.

    O ministrante nem a fomentadora responsabilizam-se por danos causados
devido a utilização das informações contidas neste material. Não há
garantias de que este material está livre de erros, assim como, todos os
sistemas em produção devem possuir {\it backup} antes de sua manipulação.

    Este material esta licenciado sobre a {\bf GNU Free Documentation
License - GFDL} ou {\bf Licença de Documentação Livre GNU} conforme descrito a
seguir:

\begin{BoxVerbatim}
    Copyright (c) 2017 Reinaldo Gil Lima de Carvalho - reinaldoc@gmail.com

    É garantida a permissão para copiar, distribuir e/ou modificar este documento
sob os termos da Licença de Documentação Livre GNU (GNU Free Documentation
License) Versão 1.3, publicada pela Free Software Foundation; com todas Seções
Secundárias Invariantes.
\end{BoxVerbatim}

    A {\bf Licença de Documentação Livre GNU} permite que todo conteúdo
esteja livre para cópia e distribuição, bem como que a propriedade autoral
seja protegida. O objetivo é garantir que o conhecimento seja livre, assim
como garante o reconhecimento ao autor.

     O autor recomenda ainda que este material seja sempre distribuído
"como está", no formato original. Contribuições e sugestões de melhorias
sobre este material podem ser enviadas ao autor e serão consideradas para
aprimoramento do material.

\section{Compilação de {\it software} X {\it software} da distribuição/sistema operacional}\setcounter{SteP}{1}

    O acesso ao código fonte do {\it software} e sua compilação, é uma das
liberdades propiciadas pelo {\it software livre}. Entretanto, o {\it software}
também pode ser obtido em forma executável (compilada), e de forma integrada
ao sistema operacional (empacotado), já estando pronto para utilização.
Cada uma destas opções possui vantagens e desvantagens que serão enumeradas
a seguir:

    Características do {\it software} obtido na forma de executável (previamente compilado):

\begin{itemize}
\item{\bf } {\bf V:} Instalação rápida que requer menos espaço em disco; evita a
compilação do {\it software}, assim como, a instalação de {\it software} de compilação
(make, gcc, etc) e cabeçalhos de bibliotecas (libc6-dev, etc)

\item{\bf } {\bf V:} Instalação automatizada de {\it software} e de bibliotecas
necessárias (dependências) para o funcionamento do {\it software} principal.

\item{\bf } {\bf V:} Versão testada pelo distribuidor do {\it software} (em geral
o distribuidor do sistema operacional), e possivelmente livre de erros.

\item{\bf } {\bf V:} Possibilita atualizações e correções de falhas de segurança
de forma automática, e fornecida pelo distribuidor do sistema operacional.

\item{\bf } {\bf V:} Facilita suporte externo devido ao método de instalação
padronizado e utilização de versões invariantes do {\it software}.

\item{\bf } {\bf V/D:} Pode não ser a versão mais nova do {\it software}, e não
possuir funcionalidades mais recentes. Todavia, a utilização de versões 
maduras, tende a fornecer maior estabilidade.

\end{itemize}


    Características do {\it software} obtido a partir do código fonte:

\begin{itemize}
\item{\bf } {\bf D:} Instalação mais complexa e demorada, demanda instalação
manual de bibliotecas externas.

\item{\bf } {\bf D:} Atualizações e correções são manuais, exigindo atenção
diária às atualizações necessárias para correções de falhas de segurança.

\item{\bf } {\bf D:} Dificulta suporte externo pois não é um método de
instalação padronizado.

\item{\bf } {\bf V/D:} Permite utilizar a última versão do {\it software},
com os novos recursos, mas trata-se de código menos testado podendo possuir
falhas não detectadas.

\item{\bf } {\bf V:} Pode permitir um ganho de performace com a compilação com
otimizações do processador, e também com o desligamento de recursos não
utilizados do {\it software}.

\end{itemize}

    Após o levantamento destas características, é notável que em ambientes
corporativos a utilização de {\it software} fornecido por um distribuidor é
essencial para continuidade da disponibilidade dos sistemas.

    Diminui-se o esforço empregado para manter o parque tecnológico atualizado
e livre de falhas. Dessa forma, o treinamento utilizará os pacotes fornecidos
pelo distribuidor do sistema operacional escolhido.

\section{Debian GNU/Linux}\setcounter{SteP}{1}

    Os sistemas operacionais desenvolvidos a partir de tecnologias livres
fornecem ferramentas que realizam instalações automatizadas de {\bf pacotes
{\it software}}. O Debian utiliza o {\bf apt-get} para este fim e fornece
repositórios web que contém pacotes de {\it software} disponíveis para
instalação. Os pacotes são contém arquivos compactados com rotinas de
pré/pós instalação e remoção, além de informações sobre dependências.

    O comportamento padrão do utilitário {\it apt-get}, ao instalar um
{\it software}, é realizar a instalação das {\bf dependências}, e também
daqueles pacotes especificados como {\bf recomendados}. Entretanto, isto
ocasiona a instalação de pacotes não requeridos e demanda a utilização de
espaço em disco adicional.

    A instalação automática de {\it software} recomendado pode ser
desabilitada por meio da adição da configuração abaixo ao arquivo
{\bf /etc/apt/apt.conf}:

\begin{BoxVerbatim}
APT::Install-Recommends "0";
\end{BoxVerbatim}

    A configuração da fonte dos pacotes a serem instalados é realizada
no arquivo {\bf /etc/apt/sources.list}:

\begin{BoxVerbatim}
deb http://ftp.br.debian.org/debian stretch main contrib non-free
deb http://security.debian.org/ stretch/updates main contrib non-free
\end{BoxVerbatim}

    Caso a conectividade seja fornecida por um proxy via http, a seguinte
configuração deve ser adicionada ao arquivo {\bf /etc/apt/apt.conf}, com a
devida adequação ao endereço IP do servidor proxy:

\begin{BoxVerbatim}
Acquire::http::Proxy "http://172.16.0.1:3128/";
\end{BoxVerbatim}

    Após a definição das fontes, é necessário o {\it download} da lista de
pacotes de {\it software} disponíveis, que é formada por informações de
versão e descrição. Esse {\it download}, bem como a atualização do sistema,
podem ser realizados pelos comandos:

\begin{BoxVerbatim}
    # apt-get update
    # apt-get upgrade
\end{BoxVerbatim}

    A lista de {\it software} disponíveis pode ser consultada, como indicado
no exemplo abaixo:

\begin{itemize}
\item{\bf }Pesquisar pelo nome do {\it software}:
\end{itemize}

\begin{BoxVerbatim}
    # apt-cache search ^apache2
\end{BoxVerbatim}

    Maiores informações sobre um determinado {\it software} podem ser obtidas
como indicado a seguir:

\begin{BoxVerbatim}
    # apt-cache show apache2
\end{BoxVerbatim}


\newpage \chapter{Instalação}\setcounter{SteP}{1}

    Este capítulo contém os procedimentos para instalação e configuração
essencial dos pacotes de software necessários para a execução do Drupal.

\section{Apache e PHP}\setcounter{SteP}{1}

\begin{BoxVerbatim}
# apt-get install apache2 libapache2-mod-php7.0
# mkdir /var/www/novoportal
\end{BoxVerbatim}

\begin{itemize}
\item{\bf }VirtualHost: /etc/apache2/sites-available/novoportal.conf
\end{itemize}

\begin{VerbatimNumerado}
<IfModule mod_ssl.c>
    <VirtualHost *:80>
        ServerName novoportal.trt8.net
        Redirect / https://novoportal.exmaple.com/
    </VirtualHost>
    <VirtualHost *:443>
        ServerAdmin reinaldoc@gmail.com
        ServerName novoportal.example.com
        DocumentRoot /var/www/novoportal
        ErrorLog ${APACHE_LOG_DIR}/error.log
        CustomLog ${APACHE_LOG_DIR}/access.log combined

        # SSL Engine
        SSLEngine on
        SSLCertificateFile    /etc/ssl/private/novoportal.cer
        SSLCertificateKeyFile /etc/ssl/private/novoportal.key

        <FilesMatch "\.php$">
            SSLOptions +StdEnvVars
        </FilesMatch>
    </VirtualHost>
</IfModule>
\end{VerbatimNumerado}

\begin{BoxVerbatim}
# a2enmod ssl
# a2ensite novoportal
\end{BoxVerbatim}

\begin{itemize}
\item{\bf }Configurações gerais
\end{itemize}

\begin{BoxVerbatim}
# /etc/apache2/conf-available/security.conf
# l.25
ServerTokens Prod
\end{BoxVerbatim}

\begin{BoxVerbatim}
# /etc/php/7.0/apache2/php.ini
# l.656
post_max_size = 20M
# l.809
upload_max_filesize = 12M
\end{BoxVerbatim}

\begin{BoxVerbatim}
# /etc/init.d/apache2 restart
\end{BoxVerbatim}

\section{MySQL}\setcounter{SteP}{1}

\begin{BoxVerbatim}
# apt-get install mariadb-server-10.1
\end{BoxVerbatim}

\begin{itemize}
\item{\bf }Preparar banco de dados
\end{itemize}

\begin{BoxVerbatim}
# mysql -p
> CREATE DATABASE novoportal CHARACTER SET utf8mb4 COLLATE utf8mb4_unicode_ci;
> CREATE USER 'novoportal'@'localhost' IDENTIFIED BY 'top-secret';
> GRANT SELECT, INSERT, UPDATE, DELETE, ALTER, CREATE, DROP, INDEX, LOCK TABLES, REFERENCES ON novoportal.* TO 'novoportal'@'localhost';
> FLUSH PRIVILEGES;
> QUIT;
\end{BoxVerbatim}

\section{FTP Server}\setcounter{SteP}{1}

\begin{BoxVerbatim}
@ instalar servidor VS FTP
# apt-get install vsftpd

@ permitir que usuários com shell /bin/false realizem login no FTP.
# echo '/bin/false' >> /etc/shells
\end{BoxVerbatim}

\begin{itemize}
\item{\bf }Configurações: /etc/vsftpd.conf
\end{itemize}

\begin{BoxVerbatim}
# l.31
write_enable=YES
# l.103
ftpd_banner=FTP Server ready.
# l.114
chroot_local_user=YES
allow_writeable_chroot=YES
# l.155
utf8_filesystem=YES
\end{BoxVerbatim}

\section{Drupal}\setcounter{SteP}{1}

\begin{itemize}
\item{\bf }https://ftp.drupal.org/files/projects/drupal-7.56.tar.gz
\item{\bf }https://ftp.drupal.org/files/translations/7.x/drupal/drupal-7.56.pt-br.po
\end{itemize}

\begin{BoxVerbatim}
adduser --home /var/www/novoportal --no-create-home --shell /bin/false novoportal
mkdir /var/www/novoportal
chown novoportal: /var/www/novoportal
su - novoportal -s /bin/bash
tar --strip-components=1 -xzf /path/to/drupal-7.xx.tar.gz
mv /path/to/drupal-7.56.pt-br.po profiles/standard/translations
\end{BoxVerbatim}

    Inicie o navegador para acesso ao Drupal.


\newpage
\bibliographystyle{unsrt}
\bibliography{Bibliography}
\newpage
\appendix
\newpage \chapter{Firewall}\setcounter{SteP}{1}

\section{Logs no console}\setcounter{SteP}{1}

\begin{itemize}
\item{\bf } {\bf /etc/sysctl.conf}: descomentar configuração a seguir (l. 10).
\end{itemize}

\begin{BoxVerbatim}
kernel.printk = 4 4 1 7
\end{BoxVerbatim}

	Aplicar configuração:

\begin{VerbatimNumerado}
sysctl -p
\end{VerbatimNumerado}

\section{Configuração}\setcounter{SteP}{1}

\begin{itemize}
\item{\bf } {\bf /etc/init.d/firewall.sh}: criar arquivo com conteúdo abaixo
e realizar adequações necessárias.
\end{itemize}

\begin{BoxVerbatim}
#!/bin/bash
# Firewall invoke script
# Copyright (c) 2012 - Reinaldo de Carvalho <reinaldoc@gmail.com>

### BEGIN INIT INFO
# Provides:          iptables
# Required-Start:    $all
# Required-Stop:     $all
# Default-Start:     2 3 4 5
# Default-Stop:      0 1 6
# Short-Description: Firewall invoke script
### END INIT INFO

case "$1" in
    start)
        iptables -F
        iptables -P INPUT DROP
        iptables -A INPUT -i lo -j ACCEPT
        iptables -A INPUT -m state --state ESTABLISHED,RELATED -j ACCEPT
        iptables -A INPUT -p tcp --syn --dport   22 -j ACCEPT -s 10.15.10.0/24
        iptables -A INPUT -p tcp --syn --dport 4848 -j ACCEPT -s 10.15.10.0/24
        iptables -A INPUT -p tcp --syn --dport 8080 -j ACCEPT
        iptables -A INPUT -p tcp --syn --dport 5666 -j ACCEPT -s 10.15.1.5
        iptables -A INPUT -p icmp --icmp-type echo-request -j ACCEPT -s 10.15.1.5
        iptables -A INPUT -m limit --limit 20/s -j LOG --log-prefix 'Firewall: '
    ;;
    stop)
        iptables -P INPUT ACCEPT
        iptables -F
    ;;
    status|list)
        iptables -nvL
    ;;
    *)
        echo "Sintaxe: $0 [ start | stop | list ]"
    ;;
esac
\end{BoxVerbatim}

    Definir permissão de execução e inicialização automática:

\begin{VerbatimNumerado}
chmod 755 /etc/init.d/firewall.sh
update-rc.d firewall.sh defaults
/etc/init.d/firewall.sh start
\end{VerbatimNumerado}

\section{FTP Conntrack}\setcounter{SteP}{1}

\begin{itemize}
\item{\bf } {\bf /etc/modules}: adicionar módulo para inicialização automática;
\end{itemize}

\begin{BoxVerbatim}
nf_conntrack_ftp
\end{BoxVerbatim}

    Esta configuração é necessária para carregar o referido módulo durante o boot,
para carréga-lo no boot atual, execute:

\begin{BoxVerbatim}
# modprobe nf_conntrack_ftp
\end{BoxVerbatim}

\newpage \chapter{Licença}\setcounter{SteP}{1}

\section{Documentação}\setcounter{SteP}{1}

\begin{BoxVerbatim}
    Copyright (c) 2017 Reinaldo Gil Lima de Carvalho - reinaldoc@gmail.com

    É garantida a permissão para copiar, distribuir e/ou modificar este documento
sob os termos da Licença de Documentação Livre GNU (GNU Free Documentation
License) Versão 1.3, publicada pela Free Software Foundation; com todas Seções
Secundárias Invariantes.
\end{BoxVerbatim}



\end{document}
